%
% nats-relay-teardown-unbounded.tex
%
% Z Specification for the force-reconnect teardown held under _connect_lock.
% Companion to nats-relay.tex, which models ForceReconnect atomically.
%
% Running `probcli nats-relay-teardown.tex -init -model_check -ctlformula
%   "AG({lock = held} => EF {lock = free})"` is TRUE.  Deleting the
% AbandonTeardown schema (the unbounded close, nats-relay-teardown-unbounded
% in the prose below) makes the same formula FALSE and the model checker
% reports a deadlock: a hung close() holds the lock forever.
%

\documentclass[a4paper,10pt,fleqn]{article}
\usepackage[margin=1in]{geometry}
\usepackage{fuzz}

\begin{document}

\title{Biff NATS Relay: Teardown Under the Connect Lock}
\date{November 2026}
\maketitle

\section{Why this is separate}

\texttt{nats-relay.tex} models $ForceReconnect$ as one atomic step because
the still-wedged case has no interesting intermediate state.  That is true
of the \emph{state} the step reaches, and false of the \emph{time} it takes:
the relay runs \texttt{safe\_close(nc)} while holding \texttt{\_connect\_lock},
and nats-py's \texttt{close()} ends in an unbounded \texttt{wait\_closed()}
that a half-open socket can stall.  While it stalls, every slow-path caller
(the poller, the heartbeat) is blocked on the lock.  This module splits the
step into begin, complete and abandon so the held-lock interval is a
modelled state, and proves the bound makes the frozen-lock state
unreachable.

\begin{zed}
  ZBOOL ::= ztrue | zfalse
\end{zed}

\begin{zed}
  LockState ::= free | held
\end{zed}

\section{State}

\begin{schema}{Teardown}
  lock : LockState \\
  inFlight : ZBOOL \\
  closeHung : ZBOOL
  \where
  inFlight = ztrue \implies lock = held \\
  closeHung = ztrue \implies inFlight = ztrue
\end{schema}

$inFlight$ is the interval in which \texttt{\_force\_reconnect} holds the
lock and awaits \texttt{safe\_close}.  $closeHung$ records that this close
will never finish by itself; it is chosen when the teardown begins, since
whether the socket is half-open is a property of the environment, not of the
relay.

\begin{schema}{Init}
  Teardown'
  \where
  lock' = free \\
  inFlight' = zfalse \\
  closeHung' = zfalse
\end{schema}

\section{Operations}

\begin{schema}{BeginTeardown}
  \Delta Teardown
  \where
  lock = free \\
  lock' = held \\
  inFlight' = ztrue
\end{schema}

The relay takes the lock and starts closing the wedged client.
$closeHung'$ is left free, so the checker explores both a close that
finishes and one that does not.

\begin{schema}{CompleteTeardown}
  \Delta Teardown
  \where
  inFlight = ztrue \\
  closeHung = zfalse \\
  lock' = free \\
  inFlight' = zfalse \\
  closeHung' = zfalse
\end{schema}

The close finished; the relay drops the client and releases the lock.

\section{Properties}

Both are checked with probcli:

\begin{quote}
  $AG\,(\,lock = held \Rightarrow EF\;lock = free\,)$
\end{quote}

\textbf{TRUE} with $AbandonTeardown$ present: from every reachable state
the lock is eventually released.  With $AbandonTeardown$ deleted
(\texttt{safe\_close} unbounded, the code before the fix) the same formula
is \textbf{FALSE}, and the model checker reports the state $lock = held \land inFlight = ztrue \land
closeHung = ztrue$ as terminal.  That state is the frozen
\texttt{\_connect\_lock}: the poller and heartbeat block on it forever.

The two modules are not yet one machine: $Teardown$ is not a component of
$Connection$ in \texttt{nats-relay.tex}, so this module proves the bound on
the lock but not its interaction with $subGen$ and $generation$.  Folding
$inFlight$ into $Connection$ is tracked as follow-up work.

\end{document}
